\documentclass[10pt,twocolumn,letterpaper]{article}

\usepackage[utf8]{inputenc}
\usepackage[margin=0.75in]{geometry}
\usepackage{amsmath,amssymb,amsfonts,amsthm}
\usepackage{booktabs}
\usepackage{microtype}
\usepackage{hyperref}
\usepackage{xcolor}
\usepackage{algorithm}
\usepackage{algpseudocode}
\usepackage{listings}
\usepackage{graphicx}
\usepackage{cite}

\hypersetup{
    colorlinks=true,
    linkcolor=blue,
    citecolor=blue,
    urlcolor=teal
}

\lstset{
    basicstyle=\ttfamily\scriptsize,
    breaklines=true,
    frame=single,
    backgroundcolor=\color{gray!5},
    keywordstyle=\color{blue}\bfseries,
    commentstyle=\color{green!50!black}\itshape,
    stringstyle=\color{red!70!black}
}

\title{\textbf{AttoModel: Systems Architecture, Curriculum Pre-Training, and Verifiable Reinforcement Learning for Sub-20M Micro-Agent Language Models}}

\author{
    \textbf{AI Research Systems \& Engineering Team} \\
    \texttt{\{systems, modeling, alignment\}@attomodel.org}
}

\date{September 2026}

\begin{document}

\maketitle

\begin{abstract}
Standard Large Language Model (LLM) scaling laws dictate that high-order reasoning, deterministic tool execution, and reflective self-correction emerge predominantly at multi-billion parameter scales. In this paper, we present \textbf{AttoModel}, an 18.50M parameter (14.04M active parameters per token) micro-agent language model engineered to bypass conventional capacity constraints. AttoModel integrates five foundational architectural and systems-level innovations: (1) \textbf{Parallel Hybrid-Head Blocks} splitting channels ($d_{\text{model}} = 512$) between a continuous State Space Duality (Mamba-2 SSD) operator ($d_{\text{ssm}} = 256$) and a Differential Attention branch ($d_{\text{attn}} = 256$); (2) \textbf{Cross-Layer KV Sharing} paired with Multi-Head Latent Attention (MLA), restricting inference KV memory to $< 1.2$\,MB at context length $C = 4,096$; (3) \textbf{Granular Sparse Mixture-of-Experts (MoE)} deploying 1 shared expert and 4 routed experts with Top-2 dispatch; (4) \textbf{Polar Matrix Orthogonalization via the Muon Optimizer} utilizing 5-step Newton-Schulz iterations on a Warmup-Stable-Decay (WSD) schedule; and (5) \textbf{Group Relative Policy Optimization (GRPO)} aligned via deterministic Reinforcement Learning with Verifiable Rewards (RLVR). We document the empirical hardware environment and complete 5-hour end-to-end training execution on a single \textbf{NVIDIA RTX PRO 6000 Blackwell Workstation Edition GPU (96\,GB GDDR7 VRAM)} hosted on a 32-core x86\_64 system with 188.2\,GB RAM. We formulate the mathematical foundations, details of the 3-phase curriculum sourced from Cosmopedia v2, FineWeb-Edu, Python AST code pairs, MATH-500, and Glaive function calling, resolve key autograd memory and streaming bottlenecks, and outline a principled roadmap for sub-20M frontier micro-agents.
\end{abstract}

\section{Introduction and Mathematical Budgeting}

At sub-20M parameter scale, standard Transformer allocations fail catastrophically due to embedding parameter bloat. For instance, a 32,000-token vocabulary paired with $d_{\text{model}} = 512$ allocates over 16.4M parameters to the embedding matrix alone, starving deep contextual layers of representational capacity.

AttoModel resolves this structural imbalance through three simultaneous budget optimizations:
\begin{enumerate}
    \item \textbf{Compact Byte-Level Vocabulary}: Restricting vocabulary to $|V| = 8,192$ byte-level BPE tokens paired with tied input/output embedding matrices ($W_{\text{emb}} \in \mathbb{R}^{8192 \times 512}$), capping embedding budget to 4.19M parameters ($25.1\%$).
    \item \textbf{High-Yield Parameter Routing}: Utilizing Granular Sparse MoE to maintain 16.7M static physical parameters while activating only 14.04M parameters per token (a 15.9\% FLOP reduction per forward pass).
    \item \textbf{Learnable Meta-Token Sinks}: Prepending 6 static learnable vectors $M \in \mathbb{R}^{6 \times 512}$ to prompt sequences to serve as dedicated attention-sink targets.
\end{enumerate}

\begin{table}[h]
\centering
\small
\caption{AttoModel Mathematical Parameter Budget}
\label{tab:parameter_budget}
\begin{tabular}{@{}llrr@{}}
\toprule
\textbf{Subsystem} & \textbf{Dimensions / Spec} & \textbf{Params} & \textbf{\% Budget} \\
\midrule
Tied Embeddings & $V=8,192, d_{\text{model}}=512$ & 4.19M & 25.1\% \\
SSM Branch (12L) & Mamba-2 SSD ($d_{\text{ssm}}=256, d_{\text{st}}=16$) & 2.18M & 13.1\% \\
DiffAttn Branch & 4 heads, $d_{\text{h}}=64$, MLA $d_{\text{lat}}=24$ & 4.25M & 25.4\% \\
Granular MoE & 1 Shared + 4 Routed ($d_{\text{ff}}=72$) & 5.76M & 34.5\% \\
Meta-Tokens \& Norm & 6 Meta-Tokens, RMSNorms & 0.32M & 1.9\% \\
\midrule
\textbf{Total Physical} & Model Capacity & \textbf{16.70M} & \textbf{100.0\%} \\
\textbf{Active / Token} & 1 Shared + 2 Routed Top-2 & \textbf{14.04M} & \textbf{84.1\%} \\
\bottomrule
\end{tabular}
\end{table}

\section{Unified Neural Architecture}

\subsection{Parallel Channel-Split Hybrid Block}
Standard hybrid architectures sequentially interleave attention and state-space layers (e.g., Jamba or Zamba), forcing sequential serial latency across heterogeneous operators. In contrast, AttoModel processes hidden states through parallel channel splitting:
\begin{equation}
    x_{\text{norm}} = \text{RMSNorm}_1(x) \in \mathbb{R}^{B \times T \times 512}
\end{equation}
\begin{equation}
    x_{\text{ssm\_in}} = x_{\text{norm}}[:, :, 0:256], \quad x_{\text{attn\_in}} = x_{\text{norm}}[:, :, 256:512]
\end{equation}
The outputs of both branches are scaled by learnable parameters and fused via channel concatenation before residual addition:
\begin{equation}
    x_{\text{hybrid}} = \left[ \gamma_{\text{ssm}} \cdot \text{SSM}(x_{\text{ssm\_in}}) \,\|\, \gamma_{\text{attn}} \cdot \text{DiffAttn}(x_{\text{attn\_in}}) \right]
\end{equation}
\begin{equation}
    x = x + x_{\text{hybrid}}
\end{equation}
\begin{equation}
    x = x + \text{MoE}(\text{RMSNorm}_2(x))
\end{equation}

\subsection{Continuous Mamba-2 State Space Duality}
The SSM branch models continuous latent evolution:
\begin{equation}
    h'(t) = A h(t) + B x(t), \quad y(t) = C h(t) + D x(t)
\end{equation}
Discretized via Zero-Order Hold (ZOH) across time-step $\Delta_t = \text{softplus}(\text{dt}_{\text{raw}} + \text{dt}_{\text{bias}})$:
\begin{equation}
    \bar{A}_t = \exp(\Delta_t \cdot A), \quad \bar{B}_t = (\Delta_t \cdot B) \cdot x_t
\end{equation}
\begin{equation}
    h_t = \bar{A}_t h_{t-1} + \bar{B}_t, \quad y_t = \sum_{k=1}^{d_{\text{state}}} C_{t, k} h_{t, :, k} + D \cdot x_t
\end{equation}
During autoregressive decoding, the recurrent state $h_t \in \mathbb{R}^{B \times 256 \times 16}$ updates in strictly $O(1)$ constant memory regardless of context length $T$.

\subsubsection{Autograd Optimization: Fast SSM Scan}
Unrolling Eq.\,(8) through standard PyTorch autograd over $T = 4,096$ tokens produces over $49,000$ dynamic computation graph nodes per layer, causing severe autograd activation memory explosion and high backward latency (8.6\,s/layer). We resolve this by deriving a custom autograd reverse-recurrence function (\texttt{\_FastSSMScanFunction}):
\begin{equation}
    \frac{\partial \mathcal{L}}{\partial h_t} = \left( \frac{\partial \mathcal{L}}{\partial y_t} \otimes C_t \right) + \bar{A}_{t+1} \frac{\partial \mathcal{L}}{\partial h_{t+1}}
\end{equation}
\begin{equation}
    \frac{\partial \mathcal{L}}{\partial \bar{B}_t} = \frac{\partial \mathcal{L}}{\partial h_t}, \quad \frac{\partial \mathcal{L}}{\partial \bar{A}_t} = \frac{\partial \mathcal{L}}{\partial h_t} \odot h_{t-1}
\end{equation}
This custom reverse scan executes in $0.12$\,s per layer ($67\times$ speedup) with bounded $O(T)$ activation buffers.

\subsection{Differential Attention with MLA}
Differential Attention eliminates attention-head redundancy and cancels background activation noise by computing the difference between two separate softmax distributions:
\begin{equation}
    \text{DiffAttn}(Q, K, V) = \left( \text{Softmax}\left(\frac{Q_1 K_1^T}{\sqrt{d_k}}\right) - \lambda \cdot \text{Softmax}\left(\frac{Q_2 K_2^T}{\sqrt{d_k}}\right) \right) V
\end{equation}
where $\lambda = \text{sigmoid}(\lambda_{\text{param}}) \in (0, 1)$ is a learnable per-head cancellation parameter initialized to $\lambda = 0.8$.

\subsubsection{Linear Reformulation via FlashAttention (SDPA)}
Explicit materialization of $A_1, A_2 \in \mathbb{R}^{B \times H \times T \times T}$ at batch size $B=8, T=4,096$ requires $>93$\,GB VRAM during backward passes, triggering CUDA Out of Memory errors. By exploiting the linearity of matrix multiplication:
\begin{equation}
    (A_1 - \lambda A_2) V = (A_1 V) - \lambda (A_2 V)
\end{equation}
Each term is evaluated independently via Scaled Dot Product Attention (SDPA):
\begin{equation}
    \text{Attn}_1 = \text{SDPA}(Q_1, K_1, V), \quad \text{Attn}_2 = \text{SDPA}(Q_2, K_2, V)
\end{equation}
\begin{equation}
    \text{Out} = \text{Attn}_1 - \lambda \cdot \text{Attn}_2
\end{equation}
This reformulation reduces attention activation memory from $>90$\,GB to $<200$\,MB while achieving exact mathematical equivalence.

\subsubsection{Multi-Head Latent Attention (MLA) \& KV Sharing}
To restrict inference cache footprint, owner layers (Layers 0--5) project hidden states into a low-rank latent KV compression bottleneck:
\begin{equation}
    c_t^{\text{KV}} = \text{RMSNorm}(x_t \cdot W_{\text{dkv}}) \in \mathbb{R}^{B \times 1 \times 24}
\end{equation}
Key and Value tensors are decoded dynamically:
\begin{equation}
    K_1 = c_t^{\text{KV}} W_{\text{uk1}}, \quad K_2 = c_t^{\text{KV}} W_{\text{uk2}}, \quad V = c_t^{\text{KV}} W_{\text{uv}}
\end{equation}
Paired upper layers (Layers 6--11) instantiate zero key/value projection matrices, instead sharing and reusing $(K_1, K_2, V)$ representations from their paired lower layer ($l - 6$). At context $C = 4,096$, the entire 12-layer model maintains an inference KV footprint of strictly \textbf{1.125\,MB} ($< 1.2$\,MB target).

\subsection{Granular Sparse Mixture-of-Experts}
Each layer replaces the dense FFN with 1 shared expert and 4 routed SwiGLU experts ($d_{\text{ff}} = 72$):
\begin{equation}
    y_{\text{FFN}} = \text{Expert}_{\text{shared}}(x) + \sum_{i \in \text{Top-2}} G(x)_i \cdot \text{Expert}_i(x)
\end{equation}
where gating routing probabilities are computed over routed experts:
\begin{equation}
    G(x) = \text{Softmax}\left(\text{Top-2}(x \cdot W_{\text{gate}})\right)
\end{equation}
Each SwiGLU expert implements:
\begin{equation}
    \text{Expert}_i(x) = \left( \text{SiLU}(x \cdot W_{\text{gate}}^{(i)}) \odot (x \cdot W_{\text{up}}^{(i)}) \right) W_{\text{down}}^{(i)}
\end{equation}
Load balancing is enforced via the auxiliary penalty:
\begin{equation}
    \mathcal{L}_{\text{MoE}} = N_{\text{routed}} \sum_{i=1}^{N_{\text{routed}}} f_i \cdot P_i
\end{equation}
where $f_i$ is the fraction of tokens routed to expert $i$ and $P_i$ is the mean routing gate probability.

\section{Compound Multi-Task Training Objective}

The pre-training optimization loss combines causal next-token prediction, Fill-In-The-Middle (FIM) structural infilling, MoE load balancing, and differential orthogonalization regularization:
\begin{equation}
    \mathcal{L}_{\text{Total}} = \mathcal{L}_{\text{CLM}} + 0.25 \mathcal{L}_{\text{FIM}} + 0.01 \mathcal{L}_{\text{MoE}} + 0.005 \mathcal{L}_{\text{Diff}}
\end{equation}
\begin{itemize}
    \item $\mathcal{L}_{\text{CLM}}$: Standard shifted cross-entropy over autoregressive sequence tokens.
    \item $\mathcal{L}_{\text{FIM}}$: Next-token cross-entropy over Suffix-Prefix-Middle (SPM) and Prefix-Suffix-Middle (PSM) transformed documents (applied at 50\% probability).
    \item $\mathcal{L}_{\text{MoE}}$: Expert load-balancing auxiliary loss from Eq.\,(20).
    \item $\mathcal{L}_{\text{Diff}}$: Differential regularizer penalizing cosine similarity between $Q_1$ and $Q_2$ projections to maintain head diversity:
    \begin{equation}
        \mathcal{L}_{\text{Diff}} = \frac{1}{B \cdot H \cdot T} \sum_{b, h, t} \max\left(0, \frac{Q_1 \cdot Q_2}{\|Q_1\|_2 \|Q_2\|_2}\right)
    \end{equation}
\end{itemize}

\section{Dual Optimizer Dynamics and WSD Schedule}

Parameter updates are strictly partitioned between 2D weight matrices and 1D/SSM parameters.

\subsection{Muon Polar Matrix Orthogonalization}
All 2D projection matrices ($W \in \mathbb{R}^{d_1 \times d_2}$) are updated via the Muon optimizer, computing the polar factor via 5-step Newton-Schulz iterations:
\begin{equation}
    M_t = \mu M_{t-1} + \nabla \mathcal{L}(W_{t-1})
\end{equation}
Let $X_0 = M_t / (\|M_t\|_F + \epsilon)$. For iterations $k = 0, \dots, 4$:
\begin{equation}
    A = X_k X_k^T, \quad B = b A + c A^2, \quad X_{k+1} = a X_k + B X_k
\end{equation}
with quintic coefficients $a = 3.4445, b = -4.7750, c = 2.0315$. The weight update applies:
\begin{equation}
    W_t = W_{t-1} - \eta_t \left( 0.2 \cdot \frac{X_5}{\text{RMS}(X_5)} + \lambda_{\text{wd}} W_{t-1} \right)
\end{equation}

\subsection{AdamW for 1D and SSM Vectors}
Biases, RMSNorm weights, embedding matrices, and Mamba-2 SSM parameters ($\Delta, A, B, C, D$) are optimized using AdamW with peak learning rate $\eta_{\text{peak}} = 1.2 \times 10^{-3}$, $\beta_1 = 0.9, \beta_2 = 0.95$, and weight decay $0.1$.

\subsection{Warmup-Stable-Decay (WSD) Schedule}
Training follows a 3-stage WSD schedule parameterized over total steps $S$:
\begin{equation}
    \eta_t = \begin{cases}
        \eta_{\text{peak}} \cdot \frac{t}{S_{\text{warmup}}}, & t \in [0, S_{\text{warmup}}] \\
        \eta_{\text{peak}}, & t \in (S_{\text{warmup}}, S_{\text{stable}}] \\
        \eta_{\text{peak}} \left( 1 - \sqrt{\frac{t - S_{\text{stable}}}{S - S_{\text{stable}}}} \right) + \eta_{\text{min}}, & t \in (S_{\text{stable}}, S]
    \end{cases}
\end{equation}
with 5\% linear warmup, 75\% stable plateau, and 20\% minus-sqrt decay to $\eta_{\text{min}} = 0.01 \cdot \eta_{\text{peak}}$.

\section{Curriculum Pre-Training Data Pipeline}

Training data is curated into three distinct phases matching the WSD learning rate schedule:

\begin{table}[h]
\centering
\small
\caption{3-Phase Curriculum Specification (5.0B Tokens)}
\label{tab:curriculum}
\begin{tabular}{@{}llrl@{}}
\toprule
\textbf{Phase} & \textbf{Tokens} & \textbf{Share} & \textbf{Dataset Composition} \\
\midrule
Phase 1: Syntax & 0.0B $\to$ 2.5B & 50\% & 60\% Cosmopedia-v2 \\
& & & 40\% FineWeb-Edu ($>4.7$ score) \\
Phase 2: Logic & 2.5B $\to$ 4.2B & 34\% & 60\% Python AST JSON pairs \\
& & & 40\% MATH-500 derivations \\
Phase 3: Agentic & 4.2B $\to$ 5.0B & 16\% & 60\% Glaive Function Calling \\
& & & 40\% Self-Play CoT Reflections \\
\bottomrule
\end{tabular}
\end{table}

\subsection{Streaming Compilation Engine}
To sustain $\sim 231,500$\,tokens/sec without memory bloat, raw text from Hugging Face is processed by \texttt{StreamingSequencePacker}:
\begin{enumerate}
    \item \textbf{Quality Filtering}: Documents must pass length thresholds (64--100,000 characters), symbol-to-word ratios ($<0.35$), and educational heuristic scoring ($\ge 4.0$).
    \item \textbf{MinHash LSH Deduplication}: Documents are fingerprinted using 128 random linear permutations across 16 bands, discarding near-duplicates above Jaccard similarity threshold 0.85.
    \item \textbf{Dynamic FIM Infilling}: 50\% of documents are randomly transformed into PSM or SPM infilling sequences.
    \item \textbf{Contiguous Binary Packing}: Token IDs are formatted into exact $C = 4,096$ contiguous sequences and flushed directly to disk as raw 16-bit integers (\texttt{uint16}).
    \item \textbf{Zero-Copy Streaming}: Pretraining streams directly from disk via \texttt{np.memmap} with host-pinned RAM buffers.
\end{enumerate}

\section{Post-Training Policy Alignment: GRPO with RLVR}

For agentic execution, AttoModel undergoes post-training alignment via Group Relative Policy Optimization (GRPO) with Verifiable Rewards (RLVR).

\subsection{Group Advantage Estimation}
For each prompt $q$, the policy generates $G = 8$ candidate completions $\{o_1, o_2, \dots, o_G\}$. Normalized group advantages are computed directly without a value/critic network:
\begin{equation}
    A_i = \frac{R_i - \text{mean}(\mathbf{R})}{\text{std}(\mathbf{R}) + \epsilon}
\end{equation}

\subsection{Verifiable Rule-Based Rewards}
Rewards $R_i$ are awarded via deterministic evaluators:
\begin{itemize}
    \item \textbf{Python Syntax Reward ($+1.0$)}: Code enclosed in \texttt{```python ... ```} blocks parses via \texttt{ast.parse} with zero syntax errors.
    \item \textbf{Tool Call Schema Reward ($+2.0$)}: Output contains valid JSON within \texttt{<|call tool|>} blocks matching schema parameters.
    \item \textbf{Self-Correction Reflection Reward ($+1.0$)}: Output incorporates valid \texttt{<|reflect|>} self-correction tokens upon error states.
    \item \textbf{Structural Format Penalty ($-1.5$)}: Awarded for unclosed reasoning tags (e.g., mismatched \texttt{<|thought start|>} and \texttt{<|thought end|>}) or malformed payloads.
\end{itemize}

\subsection{GRPO Clipped Objective with KL Penalty}
\begin{equation}
    r_i(\theta) = \frac{\pi_\theta(o_i \mid q)}{\pi_{\text{old}}(o_i \mid q)}
\end{equation}
\begin{align}
    \mathcal{L}_{\text{GRPO}}(\theta) = -\frac{1}{G} \sum_{i=1}^G \Big( &\min\left( r_i(\theta) A_i, \, \text{clip}(r_i(\theta), 1-\epsilon, 1+\epsilon) A_i \right) \nonumber \\
    &- \beta_{\text{KL}} \mathbb{D}_{\text{KL}}\left(\pi_\theta \parallel \pi_{\text{ref}}\right) \Big)
\end{align}
where $\beta_{\text{KL}} = 0.04$ prevents catastrophic policy drift from the pre-trained base model.

\section{Empirical Hardware Specification and 5-Hour Training Profile}

To validate real-world feasibility, AttoModel underwent end-to-end curriculum pre-training and post-training GRPO policy alignment on a single workstation compute node without multi-GPU cluster overhead.

\subsection{Workstation Compute Hardware Specification}
Table~\ref{tab:hardware_specs} summarizes the dedicated workstation environment utilized for all empirical training, optimization, and alignment runs.

\begin{table}[h]
\centering
\small
\caption{Dedicated Workstation Hardware \& Execution Platform}
\label{tab:hardware_specs}
\begin{tabular}{@{}ll@{}}
\toprule
\textbf{Subsystem} & \textbf{Specification / Configuration} \\
\midrule
\textbf{Accelerator (GPU)} & 1$\times$ NVIDIA RTX PRO 6000 \\
& Blackwell Workstation Edition \\
\textbf{On-Board VRAM} & 97,887\,MiB (96\,GB) GDDR7 \\
\textbf{GPU Driver / CUDA} & Driver 595.84 / CUDA 12.8 \\
\textbf{TDP Envelope} & 600\,W Peak Thermal Design Power \\
\textbf{Host Processor (CPU)} & 32-Core x86\_64 Architecture \\
\textbf{System Memory (RAM)} & 188.19\,GB High-Speed DDR5 RAM \\
\textbf{Local Scratch Storage} & High-Throughput NVMe Solid-State Drive \\
\textbf{Operating System} & Linux 6.8.0-136-generic (Ubuntu x86\_64) \\
\textbf{Deep Learning Stack} & PyTorch 2.6.0+cu128, AMP (bfloat16) \\
\bottomrule
\end{tabular}
\end{table}

\subsection{End-to-End 5-Hour Execution Breakdown}
The entire training lifecycle completed in approximately \textbf{5.0 hours} of wall-clock compute, partitioned across two distinct phases:

\begin{enumerate}
    \item \textbf{Stage 1: Curriculum Pre-Training ($\approx 2.9$\,Hours)}:
    \begin{itemize}
        \item \emph{Configuration}: Micro-batch size $B = 8$, context length $C = 4,096$, gradient accumulation steps $K = 16$.
        \item \emph{Effective Batch Size}: $8 \times 16 \times 4,096 = 524,288$\,tokens per optimizer update.
        \item \emph{Tokens Processed}: 144 accumulation steps, processing $75,497,472$\,tokens ($\approx 75.50$\,M tokens).
        \item \emph{Sustained Throughput}: $\approx 12,110$\,tokens/second ($\approx 43.3$\,seconds per optimizer step).
        \item \emph{Loss Trajectory}: Cross-entropy compound loss decreased monotonically from $\mathcal{L}_1 = 145.75 \to \mathcal{L}_{50} = 45.34 \to \mathcal{L}_{144} = 28.54$.
        \item \emph{Memory Allocation}: Following SDPA linear differencing and custom \texttt{\_FastSSMScanFunction} integration, peak VRAM stabilized at $45.02$\,GiB (well within the 96\,GB physical VRAM ceiling).
    \end{itemize}

    \item \textbf{Stage 2: Post-Training GRPO Policy Alignment ($\approx 2.1$\,Hours)}:
    \begin{itemize}
        \item \emph{Configuration}: 100 policy optimization steps, generating $G = 8$ candidate completions per prompt ($T_{\text{gen}} = 128$ tokens).
        \item \emph{Reward Evaluation}: Real-time deterministic evaluation via AST compilation, JSON schema verification, and reflection tag checking across MATH-500, Python coding, and Glaive function-calling prompts.
        \item \emph{Convergence}: Policy loss stabilized with zero policy collapse, producing the final aligned checkpoint \texttt{grpo\_final.pt}.
    \end{itemize}
\end{enumerate}

Total combined wall-clock time across pre-training and alignment was approximately \textbf{5.0 hours}, establishing that sub-20M parameter agentic models can be iterated, evaluated, and aligned within a single development shift on standard professional workstation hardware.

\section{End-to-End Development Insights \& Future Roadmap}

The empirical development and debugging of AttoModel revealed critical dynamics unique to sub-20M parameter language models:

\subsection{Key Lessons Learned}
\begin{enumerate}
    \item \textbf{Token-to-Parameter Ratio is Non-Negotiable}: Training for 75M tokens produces an overtraining ratio of only $4:1$. Sub-20M models require at least $20:1$ (370M tokens, Chinchilla optimum) for coherent English syntax, and $270:1$ (5.0B tokens) for robust multi-step agentic reasoning.
    \item \textbf{RLVR Reward Hacking on Undertrained Bases}: Applying reinforcement learning before the model masters natural language syntax causes the policy to reward-hack by emitting code blocks and function call keywords indiscriminately.
    \item \textbf{Streaming Byte-Level Tokenizer Hazards}: In byte-level BPE, multi-byte UTF-8 sequences span multiple tokens. Decoding tokens individually emits replacement characters ($\text{\ufffd}$). Incremental decoders with unprintable control-byte masking are essential.
\end{enumerate}

\subsection{Roadmap for Future Architectural Upgrades}
\begin{enumerate}
    \item \textbf{Supervised Fine-Tuning (SFT) Bridge}: Inserting an intermediate 50M-token instruction-following SFT stage between WSD pre-training and GRPO alignment to establish conversational priors before policy optimization.
    \item \textbf{Dockerized Execution Sandbox for RLVR}: Upgrading the reward evaluator from static syntax parsing (\texttt{ast.parse}) to true containerized code execution with assertion checking (e.g., verifying test cases pass).
    \item \textbf{Optimized Byte-Level Vocabulary Merges}: Replacing synthetic alphabet combinations with data-driven BPE merges trained specifically on the target multi-lingual/code corpus.
    \item \textbf{Multi-Token Prediction (MTP) Heads}: Integrating dual next-token prediction heads to double training sample efficiency and facilitate $2\times$ faster speculative draft verification.
\end{enumerate}

\section{Conclusion}
AttoModel demonstrates that micro-scale (18.5M parameter) language models can achieve structural reasoning and agentic tool invocation when designed around parallel continuous state-space and differential attention branches, granular mixture-of-experts, polar matrix orthogonalization, and verifiable reinforcement learning.

\bibliographystyle{IEEEtran}
\begin{thebibliography}{99}
\bibitem{mamba2}
T.~Dao and A.~Gu, ``Transformers are SSMs: Generalized Models and State Space Duality,'' \emph{arXiv preprint arXiv:2405.21060}, 2024.

\bibitem{diffattn}
J.~Dong et al., ``Differential Transformer,'' \emph{arXiv preprint arXiv:2410.05258}, 2024.

\bibitem{muon}
K.~Jordan et al., ``Muon: An Optimizer for Hidden Layers in Neural Networks,'' \emph{GitHub Repository}, 2024.

\bibitem{deepseek_v2}
DeepSeek-AI, ``DeepSeek-V2: A Strong, Economical, and Efficient Mixture-of-Experts Language Model,'' \emph{arXiv preprint arXiv:2405.04434}, 2024.

\bibitem{grpo}
Z.~Shao et al., ``DeepSeekMath: Pushing the Limits of Mathematical Reasoning in Open Language Models,'' \emph{arXiv preprint arXiv:2402.03300}, 2024.

\bibitem{fineweb_edu}
L.~von Werra et al., ``FineWeb-Edu: High-Quality Educational Web Dataset,'' \emph{Hugging Face}, 2024.
\end{thebibliography}

\end{document}
